% Riwayat singkat penulis, wajib menurut Permendikbudristek 22/2022 Lampiran
% III Tabel 3. Datanya sama dengan riwayat penulis pada buku saudara di
% repositori python-programming-fundamentals dan advanced-software-engineering.
\chapter{Riwayat Penulis}

Alfa Yohannis adalah dosen Universitas Pradita di Tangerang. Gelar Doctor of
Philosophy bidang ilmu komputer diperolehnya dari University of York. Gelar
Magister Teknik bidang Informatika diperoleh dari Institut Teknologi Bandung,
sedangkan gelar Sarjana Teknik bidang Teknik Elektro dari Universitas
Diponegoro.

Sebelumnya, Alfa Yohannis menjadi \textit{research fellow} di University of York,
\textit{Knowledge Transfer Partnership} (KTP) \textit{associate} di Leonardo
Ltd, dan dosen di Kalbis Institute.

Bidang keahliannya meliputi rekayasa perangkat lunak berbasis model, sistem
informasi, gamifikasi, sistem terdesentralisasi, dan pendidikan ilmu komputer.
Di Universitas Pradita, Alfa Yohannis mengampu mata kuliah Arsitektur Perangkat
Lunak (IF231303), yang bahannya menjadi dasar buku ini.

Karya terbitannya antara lain \textit{Turning Models Inside Out} (2017),
\textit{Towards Model-Based Bias Mitigation in Machine Learning} (2022), dan
\textit{LLM-Assisted TOGAF Preliminary Phase for Student Digital Twins} (2025).
Penelitian terakhirnya mengukur jejak karbon dan kinerja penyimpanan data citra
deret waktu pada PostgreSQL, MongoDB, dan MinIO.
